% !TEX root = ../dabirestan.tex
%% ============================================================
%%  فصل ۳ — احتمال در فضاهای هم‌شانس و قوانین احتمال
%% ============================================================
\chapter{احتمال در فضاهای هم‌شانس و قوانین احتمال}
\label{hs:ch3}

\chapterintro{\begin{itemize}
  \item وقتی برآمدها هم‌شانس‌اند، احتمال را با فرمول $\P(A)=\frac{n(A)}{n(S)}$ حساب کنید؛
  \item با ابزارهای شمارش فصل ۲ مسئله‌های احتمال (تاس، کارت، مهره) را حل کنید؛
  \item قانون متمم و قانون جمع را بفهمید و به کار ببرید.
\end{itemize}}

%% ------------------------------------------------------------
\section{احتمال کلاسیک}
%% ------------------------------------------------------------

اگر تاسی سالم باشد، هیچ وجهی بر وجه دیگر برتری ندارد؛ می‌گوییم برآمدها \term{هم‌شانس}{equally likely}اند. در این حالت احتمال را بدون آزمایش و فقط با شمردن حساب می‌کنیم.

\begin{definition}[label=hs:def:classical]{احتمال در فضای هم‌شانس}
اگر همهٔ برآمدهای فضای نمونهٔ $S$ هم‌شانس باشند، احتمال پیشامد $A$ برابر است با
\[ \P(A)=\frac{n(A)}{n(S)}=\frac{\text{\rl{تعداد حالت‌های مطلوب}}}{\text{\rl{تعداد کل حالت‌ها}}}. \]
\end{definition}

\begin{example}[label=hs:ex:die]{تاس}
تاس سالمی پرتاب می‌شود. احتمال پیشامدهای «زوج آمدن» و «آمدن عددی بیشتر از ۴» را بیابید.
\begin{solution}
$n(S)=6$. \ زوج: $A=\{2,4,6\}$، پس $\P(A)=\tfrac36=\tfrac12$. \ بیشتر از ۴: $B=\{5,6\}$، پس $\P(B)=\tfrac26=\tfrac13$.
\end{solution}
\end{example}

\begin{example}[label=hs:ex:twocoins]{دو سکه}
دو سکهٔ سالم پرتاب می‌شوند. احتمال اینکه «یکی شیر و یکی خط» بیاید چقدر است؟
\begin{solution}
$S=\{HH,HT,TH,TT\}$ با ۴ برآمد هم‌شانس. پیشامد مطلوب $\{HT,TH\}$ دو عضو دارد؛ پس احتمال $\tfrac24=\tfrac12$ است. (اگر به اشتباه فقط سه حالت «دو شیر، دو خط، یکی از هر کدام» را در نظر می‌گرفتیم، جواب غلطِ $\tfrac13$ به دست می‌آمد؛ آن سه حالت هم‌شانس نیستند.)
\end{solution}
\end{example}

\begin{example}[label=hs:ex:twodice]{مجموع دو تاس}
دو تاس سالم (یکی سفید و یکی قرمز) پرتاب می‌شوند. احتمال اینکه مجموع دو عدد ۷ شود چقدر است؟ مجموع ۱۲ چطور؟
\begin{solution}
طبق اصل ضرب $n(S)=6\times6=36$ برآمد هم‌شانس داریم. جدول زیر مجموع هر برآمد را نشان می‌دهد:
\begin{figblock}
\begin{tikzpicture}[x=0.82cm, y=0.62cm]
  \foreach \i in {1,...,6} \foreach \j in {1,...,6} {
    \pgfmathtruncatemacro{\s}{\i+\j}
    \ifnum\s=7 \fill[accent!40] (\i-0.5,-\j+0.5) rectangle (\i+0.5,-\j-0.5); \fi
    \node[font=\small] at (\i,-\j) {\lr{\s}};
  }
  \draw[primary, thick] (0.5,-0.5) rectangle (6.5,-6.5);
  \foreach \i in {1,...,6} { \node[font=\small\bfseries, text=defcolor] at (\i,0.3) {\lr{\i}}; \node[font=\small\bfseries, text=warncolor] at (7.2,-\i) {\lr{\i}}; }
  \node[font=\scriptsize, text=defcolor] at (3.5,1.1) {\rl{تاس سفید}};
  \node[font=\scriptsize, text=warncolor] at (8.4,-3.5) {\rl{تاس قرمز}};
\end{tikzpicture}

\figcaption[hs:fig:dice-table]{جدول مجموع دو تاس؛ خانه‌های رنگی مجموع ۷ را نشان می‌دهند.}
\end{figblock}
مجموع ۷ در ۶ خانه آمده است: $\P(\fa{مجموع ۷})=\tfrac{6}{36}=\tfrac16$. \ مجموع ۱۲ فقط یک خانه دارد: $\P(\fa{مجموع ۱۲})=\tfrac1{36}$.
\end{solution}
\end{example}

\begin{warning}
فرمول $\frac{n(A)}{n(S)}$ \emph{فقط} وقتی درست است که برآمدها هم‌شانس باشند. مثلاً اگر فضای نمونهٔ مجموع دو تاس را $\{2,3,\dots,12\}$ بگیرید (۱۱ حالت)، این حالت‌ها هم‌شانس نیستند و نمی‌توانید بگویید احتمال مجموع ۷ برابر $\tfrac1{11}$ است.
\end{warning}

\subsection{مسائل با کارت و مهره}

یک دست کارت بازی ۵۲ کارت دارد: ۴ خال (\lr{$\spadesuit$} پیک، \lr{$\heartsuit$} دل، \lr{$\diamondsuit$} خشت، \lr{$\clubsuit$} گشنیز) و در هر خال ۱۳ کارت (آس، ۲ تا ۱۰، سرباز، بی‌بی، شاه).

\begin{example}[label=hs:ex:card]{یک کارت}
یک کارت به تصادف انتخاب می‌کنیم. احتمال اینکه (الف) آس باشد؛ (ب) دل باشد؛ (ج) شاهِ دل باشد چقدر است؟
\begin{solution}
$n(S)=52$. \ (الف) ۴ آس: $\tfrac4{52}=\tfrac1{13}$. \ (ب) ۱۳ دل: $\tfrac{13}{52}=\tfrac14$. \ (ج) فقط یک کارت: $\tfrac1{52}$.
\end{solution}
\end{example}

\begin{example}[label=hs:ex:urn]{دو مهره از جعبه}
جعبه‌ای ۵ مهرهٔ قرمز و ۳ مهرهٔ آبی دارد. ۲ مهره به تصادف (با هم) برمی‌داریم. احتمال اینکه (الف) هر دو قرمز باشند؛ (ب) یکی قرمز و یکی آبی باشد چقدر است؟
\begin{solution}
ترتیب مهم نیست؛ $n(S)=\binom82=28$.\\
(الف) ۲ قرمز از ۵: $\binom52=10$؛ پس $\P=\tfrac{10}{28}=\tfrac5{14}$.\\
(ب) ۱ قرمز از ۵ و ۱ آبی از ۳: $5\times3=15$؛ پس $\P=\tfrac{15}{28}$.
\end{solution}
\end{example}

\begin{example}[label=hs:ex:class]{انتخاب نماینده}
کلاسی ۱۲ دختر و ۸ پسر دارد. ۳ نفر به تصادف برای یک اردو انتخاب می‌شوند. احتمال اینکه هر سه دختر باشند چقدر است؟
\begin{solution}
$\P=\dfrac{\binom{12}3}{\binom{20}3}=\dfrac{220}{1140}=\dfrac{11}{57}\approx0.19$.
\end{solution}
\end{example}

%% ------------------------------------------------------------
\section{قوانین احتمال}
%% ------------------------------------------------------------

از فرمول $\P(A)=\frac{n(A)}{n(S)}$ چند قانون ساده ولی بسیار مهم به دست می‌آید. این قانون‌ها در حالت کلی (حتی وقتی برآمدها هم‌شانس نیستند) هم درست‌اند.

\begin{theorem}[label=hs:thm:basic]{قوانین پایهٔ احتمال}
برای هر پیشامد $A$:
\begin{enumerate}
  \item $0\le\P(A)\le1$؛
  \item $\P(S)=1$ (پیشامد حتمی) و $\P(\varnothing)=0$ (پیشامد غیرممکن)؛
  \item \textbf{قانون متمم:} \ $\P(A')=1-\P(A)$.
\end{enumerate}
\end{theorem}

\textbf{چرا قانون متمم درست است؟} هر برآمد یا در $A$ است یا در $A'$؛ پس $n(A)+n(A')=n(S)$. با تقسیم بر $n(S)$ داریم $\P(A)+\P(A')=1$.

\begin{example}[label=hs:ex:atleast]{دست‌کم یک شیر}
سه سکهٔ سالم پرتاب می‌شوند. احتمال اینکه \emph{دست‌کم یک} شیر بیاید چقدر است؟
\begin{solution}
$n(S)=2^3=8$. متمم «دست‌کم یک شیر» این است: «هیچ شیری نیاید»، یعنی فقط $TTT$. پس
\[ \P(\fa{دست‌کم یک شیر})=1-\P(TTT)=1-\tfrac18=\tfrac78. \]
\end{solution}
\end{example}

\begin{note}
هر وقت در صورت مسئله «\textbf{دست‌کم یکی}» دیدید، اول به متمم فکر کنید: «\textbf{هیچ‌کدام}» معمولاً خیلی راحت‌تر حساب می‌شود.
\end{note}

\begin{theorem}[label=hs:thm:add]{قانون جمع}
\begin{itemize}
  \item اگر $A$ و $B$ \textbf{ناسازگار} باشند: \quad $\P(A\cup B)=\P(A)+\P(B)$.
  \item در حالت کلی: \quad $\P(A\cup B)=\P(A)+\P(B)-\P(A\cap B)$.
\end{itemize}
\end{theorem}

\textbf{چرا؟} در $\P(A)+\P(B)$، برآمدهایی که هم در $A$ و هم در $B$ هستند (ناحیهٔ مشترک در نمودار ون) \emph{دو بار} شمرده شده‌اند؛ پس یک بار $\P(A\cap B)$ را کم می‌کنیم. اگر ناسازگار باشند، ناحیهٔ مشترکی نیست.

\begin{example}[label=hs:ex:acehearts]{آس یا دل}
یک کارت به تصادف انتخاب می‌کنیم. احتمال اینکه آس \emph{یا} دل باشد چقدر است؟
\begin{solution}
$\P(\fa{آس})=\tfrac4{52}$ و $\P(\fa{دل})=\tfrac{13}{52}$. ولی آسِ دل در هر دو شمرده شده است: $\P(\fa{آس دل})=\tfrac1{52}$. پس
\[ \P(\fa{آس یا دل})=\frac4{52}+\frac{13}{52}-\frac1{52}=\frac{16}{52}=\frac4{13}. \]
\end{solution}
\end{example}

\begin{example}[label=hs:ex:sports]{ورزش دانش‌آموزان}
در یک کلاس ۳۰ نفره، ۱۸ نفر فوتبال، ۱۲ نفر والیبال و ۷ نفر هر دو ورزش را بازی می‌کنند. اگر یک نفر را به تصادف انتخاب کنیم، احتمال اینکه (الف) دست‌کم یکی از این دو ورزش را بازی کند؛ (ب) هیچ‌کدام را بازی نکند چقدر است؟
\begin{solution}
(الف) $\P(F\cup V)=\tfrac{18}{30}+\tfrac{12}{30}-\tfrac{7}{30}=\tfrac{23}{30}$.
\ (ب) طبق قانون متمم: $1-\tfrac{23}{30}=\tfrac7{30}$.
نمودار ون این را نشان می‌دهد:
\begin{figblock}
\begin{tikzpicture}[scale=0.9]
  \draw[graytext, thick, rounded corners=5pt] (-2.4,-1.7) rectangle (4.2,1.7);
  \fill[setA, opacity=0.25] (0,0) circle (1.3); \fill[setB, opacity=0.25] (1.8,0) circle (1.3);
  \draw[setA!80!black, thick] (0,0) circle (1.3); \draw[setB!80!black, thick] (1.8,0) circle (1.3);
  \node[font=\small] at (-0.45,0) {\lr{11}}; \node[font=\small] at (0.9,0) {\lr{7}}; \node[font=\small] at (2.25,0) {\lr{5}};
  \node[font=\small] at (3.75,-1.35) {\lr{7}};
  \node[font=\small, text=setA!80!black] at (-0.9,1.55) {\rl{فوتبال}};
  \node[font=\small, text=setB!80!black] at (2.7,1.55) {\rl{والیبال}};
\end{tikzpicture}

\figcaption[hs:fig:sports]{۱۱ نفر فقط فوتبال، ۷ نفر هر دو، ۵ نفر فقط والیبال و ۷ نفر هیچ‌کدام.}
\end{figblock}
\end{solution}
\end{example}

\begin{moreinfo}{مسئلهٔ روز تولد}
در یک کلاس ۲۳ نفره، احتمال اینکه دست‌کم دو نفر روز تولد یکسان داشته باشند چقدر است؟ حدس بزنید! جواب شگفت‌انگیز است: \textbf{بیش از ۵۰٪}. با قانون متمم: احتمال اینکه همه روز تولد متفاوت داشته باشند
\[ \frac{365}{365}\times\frac{364}{365}\times\frac{363}{365}\times\cdots\times\frac{343}{365}\approx0.493 \]
است؛ پس احتمال دست‌کم یک تکرار حدود $1-0.493=0.507$ است. با ۵۰ نفر این احتمال به حدود ۹۷٪ می‌رسد!
\end{moreinfo}

%% ------------------------------------------------------------
\section*{خلاصهٔ فصل ۳}
\addcontentsline{toc}{section}{خلاصه و تمرین‌ها}
%% ------------------------------------------------------------

\begin{summary}[title={\leavevmode\boxicon{accent}{$\Sigma$}\ \ خلاصهٔ فصل}]
\begin{itemize}
  \item برای برآمدهای \textbf{هم‌شانس}: $\P(A)=\frac{n(A)}{n(S)}$. قبل از استفاده مطمئن شوید برآمدها واقعاً هم‌شانس‌اند.
  \item $0\le\P(A)\le1$، \ $\P(S)=1$، \ $\P(\varnothing)=0$.
  \item \textbf{قانون متمم:} $\P(A')=1-\P(A)$؛ بهترین ابزار برای «دست‌کم یکی».
  \item \textbf{قانون جمع:} $\P(A\cup B)=\P(A)+\P(B)-\P(A\cap B)$؛ برای پیشامدهای ناسازگار جملهٔ آخر صفر است.
\end{itemize}
\end{summary}

\subsection*{تمرین‌ها}

\begin{exercise}
دو تاس سالم پرتاب می‌شوند. احتمال اینکه (الف) مجموع ۸ شود؛ (ب) دو عدد برابر بیایند (جفت)؛ (ج) مجموع کمتر از ۵ شود چقدر است؟
\answer{(الف) $\tfrac5{36}$ \quad (ب) $\tfrac16$ \quad (ج) $\tfrac16$}
\end{exercise}

\begin{exercise}
یک کارت از دست ۵۲ کارتی انتخاب می‌شود. احتمال اینکه (الف) کارت تصویری (سرباز، بی‌بی یا شاه) باشد؛ (ب) قرمز (دل یا خشت) یا شاه باشد چقدر است؟
\answer{(الف) $\tfrac{12}{52}=\tfrac3{13}$ \quad (ب) $\tfrac{26+4-2}{52}=\tfrac7{13}$}
\end{exercise}

\begin{exercise}
از اعداد ۱ تا ۳۰ یک عدد به تصادف انتخاب می‌شود. احتمال اینکه بر ۲ \emph{یا} بر ۳ بخش‌پذیر باشد چقدر است؟
\answer{$\tfrac{15+10-5}{30}=\tfrac23$}
\end{exercise}

\begin{exercise}
جعبه‌ای ۴ مهرهٔ سفید و ۶ مهرهٔ سیاه دارد. ۳ مهره به تصادف برمی‌داریم. احتمال اینکه (الف) هر سه سیاه باشند؛ (ب) دست‌کم یک سفید باشد چقدر است؟
\answer{(الف) $\tfrac{20}{120}=\tfrac16$ \quad (ب) $\tfrac56$}
\end{exercise}

\begin{exercise}
اگر $\P(A)=0.5$، $\P(B)=0.4$ و $\P(A\cap B)=0.2$، مقادیر $\P(A\cup B)$ و $\P\bigl((A\cup B)'\bigr)$ را بیابید.
\answer{$0.7$ و $0.3$}
\end{exercise}

\begin{exercise}
چهار سکهٔ سالم پرتاب می‌شوند. احتمال اینکه دقیقاً ۲ شیر بیاید چقدر است؟ دست‌کم یک خط چطور؟
\answer{$\tfrac{\binom42}{16}=\tfrac6{16}=\tfrac38$؛ \ $1-\tfrac1{16}=\tfrac{15}{16}$}
\end{exercise}
